\textit{
 Crowd density monitoring in the Malioboro tourist area, which accommodates 
 millions of visitors each year, still relies on manual observation of CCTV 
 footage, while standard detection datasets such as COCO do not cover 
 vehicles unique to the area, such as \textit{andong}, \textit{becak}, and 
 \textit{bajaj}, leaving a gap that this study addresses by building a 
 dedicated local dataset, selecting a suitable detection model, and 
 deploying it within a working monitoring system. A dataset of 1,581 
 annotated images covering nine object classes was constructed from footage 
 of 22 public CCTV cameras using semi-automatic annotation assisted by a 
 pretrained model with manual correction, after which eight variants of 
 YOLOv8 and YOLO11 were trained and evaluated on a held-out test set using 
 mAP50, mAP50-95, and F1-score, with the best variant further optimized 
 through input resolution tuning before being integrated into a web-based 
 system that processes live CCTV streams and issues density alerts based 
 on the pedestrian level of service concept. The resulting dataset 
 exhibited a severe class imbalance, with a 237-to-1 ratio between the 
 most and least frequent classes. Among the eight variants, yolov8l 
 achieved the highest mAP50 and mean F1 (0.6175 and 0.6222), while yolo11l 
 achieved the highest mAP50-95 (0.4012) at a model size 42\% smaller, 
 leading to its selection despite not leading on every metric. Tuning its 
 input resolution from 640 to 960 pixels then raised its mean F1 from 
 0.5962 to 0.6739, surpassing every variant at the base resolution. On the 
 GTX 1650 deployment device, the selected model also ran faster than 
 yolov8l (9.73 versus 8.80 frames per second), and the implemented system 
 correctly displayed live detection, density charts, and alert history 
 during testing. These findings indicate that a deliberately balanced 
 choice of architecture, model size, and resolution tuning, rather than 
 the single highest-scoring variant, yields a detection model that is well 
 suited for real-time deployment on resource-constrained hardware, although 
 the underlying class imbalance continues to limit detection performance 
 on minority classes such as bicycles and \textit{bajaj}.
}

\noindent\textbf{Keywords} : object detection, YOLO, visitor density, class imbalance, Malioboro